\documentclass[12pt]{article}

\usepackage[margin=1in]{geometry}
\usepackage[T1]{fontenc}
\usepackage{times}
\usepackage{amsmath}
\usepackage{xcolor}
\usepackage{enumitem}
\usepackage{booktabs}
\usepackage{tabularx}
\usepackage{longtable}
\usepackage{graphicx}
\usepackage{listings}
\usepackage[hidelinks]{hyperref}
\usepackage{titlesec}

\setlength{\parskip}{0.5em}
\setlength{\parindent}{0pt}
\linespread{1.0}

\titleformat{\section}{\large\bfseries}{\thesection}{0.6em}{}
\titleformat{\subsection}{\normalsize\bfseries}{\thesubsection}{0.6em}{}

\usepackage{mdframed}
\newcommand{\todo}[1]{\textcolor{red}{\textbf{[TODO: #1]}}}
\newenvironment{todobox}%
  {\begin{mdframed}[backgroundcolor=yellow!20,linecolor=red!60,linewidth=1pt,
     innertopmargin=6pt,innerbottommargin=6pt]\color{red!80!black}\small\textbf{TODO --- action required before submission.}\par\vspace{0.3em}\normalcolor}%
  {\end{mdframed}}
\sloppy

\lstset{
  basicstyle=\ttfamily\footnotesize,
  breaklines=true,
  frame=single,
  columns=fullflexible,
  keepspaces=true,
}

\begin{document}

% ======================================================================
% COVER PAGE
% ======================================================================
\thispagestyle{empty}
\vspace*{1.2in}
\begin{center}
  {\LARGE\bfseries Identifying Coordinated Amazon Review Groups\\[0.3em]
   with Distributed Text and Graph Analysis\par}

  \vspace{1.0in}

  {\large CSC 7740 --- Big Data Analytics\\
   Final Project Report\par}

  \vspace{0.5in}

  {\large Submission Date: \todo{submission date}\par}

  \vspace{0.8in}

  {\large\bfseries Team Members\par}
  \vspace{0.4em}
  \begin{tabular}{ll}
    Souhardya Saha Dip    & \texttt{ssahad1@lsu.edu} \\
    Riyadil Zannat        & \texttt{rzanna2@lsu.edu} \\
    Asif Faisal Chowdhury & \texttt{achowd6@lsu.edu} \\
  \end{tabular}

  \vspace{0.8in}

  \begin{minipage}{0.85\textwidth}
    \centering
    \textbf{Prior CSC 4740 Enrollment}\\[0.3em]
    No member of this team has previously taken CSC 4740 at LSU.
  \end{minipage}
\end{center}

\newpage

% ======================================================================
% TABLE OF CONTENTS
% ======================================================================
\thispagestyle{empty}
\tableofcontents

\newpage
\setcounter{page}{1}

% ======================================================================
\section{Introduction}
% ======================================================================

\subsection{Motivation}

Product reviews are the main signal shoppers use when they cannot inspect an
item directly, and that makes them worth manipulating. The form of
manipulation that matters at scale is not the single fake review but the
\emph{coordinated group}: a set of accounts that move together, reviewing the
same products inside the same narrow time windows, often with similar text and
similar ratings. A single review in isolation looks unremarkable. The
coordination only becomes visible when many accounts are compared against each
other across the entire corpus.

That comparison is inherently a big-data problem. Deciding whether two accounts
co-occur repeatedly requires grouping every review by product and time window,
then examining every pair of accounts inside every group. On the four Amazon
categories we study, that means 24.4 million reviews collapsing into 20.6
million candidate product-time groups, from which account pairs must be
generated and counted. The pair-generation step alone is quadratic in group
size, so the work cannot be done by reading the data once on a single machine
with a hash map; it needs a shuffle-based distributed engine that can spill to
disk and parallelise the self-join.

The practical significance is twofold. For a platform, a ranked list of
suspicious account groups is a triage queue for human moderators. For
researchers, the same pipeline is a measurement instrument: it quantifies how
much apparent coordination exists in a public corpus, and how sensitive that
quantity is to the thresholds an analyst chooses.

\subsection{Project Description}

We built a distributed Spark pipeline that ingests the raw Amazon Reviews'23
corpus, enriches it with product metadata, constructs an account co-review
graph, extracts connected components as candidate coordinated groups, and
ranks those groups by a seven-signal coordination score.

The pipeline has five stages:

\begin{enumerate}
  \item \textbf{Ingest and enrich.} Read gzip-compressed JSONL reviews under an
        explicit schema, join to product metadata on \texttt{parent\_asin},
        and derive per-review features including deviation from the product's
        catalogue average rating.
  \item \textbf{Candidate grouping.} Bucket each review into a
        (product, time-window) key. Windows are 24 hours by default, narrowed
        to 6 hours for products with more than 100 reviews in a window so that
        popular items do not generate spurious co-occurrence.
  \item \textbf{Graph construction.} Self-join accounts within each group to
        produce co-occurring pairs, keep only pairs appearing together in at
        least three distinct groups, and audit the survivors for
        same-underlying-account artefacts.
  \item \textbf{Community detection.} Run distributed connected components
        over the pair graph using GraphFrames, with a Pregel-style label
        propagation implementation as a dependency-free fallback.
  \item \textbf{Scoring and output.} Compute per-component statistics,
        including MinHash LSH text similarity, combine them into a single
        coordination score, and write ranked results to Parquet.
\end{enumerate}

Beyond the detector itself, we treated the project as a measurement study. We
ran a parameter sensitivity sweep across time windows and co-occurrence
thresholds, validated recall against synthetically planted coordination rings,
and measured scale-out behaviour by executing the identical pipeline on a
one-worker and a two-worker AWS EMR cluster.

\subsection{Division of Roles}

\begin{todobox}
The team must confirm and correct this table before submission. It is filled
in from the repository's commit history and the division of work as executed,
but each member should verify their own row.
\end{todobox}

\begin{tabularx}{\textwidth}{@{}p{1.5in}>{\raggedright\arraybackslash}X@{}}
  \toprule
  \textbf{Member} & \textbf{Contributions and Responsibilities} \\
  \midrule
  Souhardya Saha Dip &
    Pipeline architecture and the DataFrame implementation
    (\texttt{AmazonReviewCoordinationDF.py}); schema definitions and the
    metadata join; the seven-signal scoring function; GraphFrames integration
    and the label-propagation fallback; same-account pair auditing; AWS EMR
    cluster provisioning, S3 data staging, and the one-worker versus
    two-worker scale-out experiment; parameter sensitivity sweep and
    synthetic injection validation; performance profiling and this report. \\
  \addlinespace
  Riyadil Zannat &
    \todo{confirm} Repository setup and the containerised development
    environment (\texttt{Dockerfile}, \texttt{.devcontainer}); initial
    pipeline prototype and per-category run scripts
    (\texttt{run\_all\_categories.sh}, \texttt{run\_local.sh}); dataset
    acquisition scripting. \\
  \addlinespace
  Asif Faisal Chowdhury &
    \todo{confirm} Exploratory analysis scripts
    (\texttt{analysis/feasibility.py}, \texttt{analysis/pairest.py},
    \texttt{analysis/lockstep.py}, \texttt{analysis/nearDup.py}) used to size
    the problem and choose the co-occurrence threshold; sample dataset
    construction for the deliverable package. \\
  \bottomrule
\end{tabularx}

% ======================================================================
\section{System Architecture and Design}
% ======================================================================

\subsection{Chosen Big Data Frameworks}

\textbf{Apache Spark (PySpark DataFrame API), version 3.5.} Spark is the right
engine here because the dominant cost is a \emph{shuffle}, not a scan. Pair
generation requires regrouping 24.4 million reviews by a synthetic
(product, time-window) key and then self-joining within each group; community
detection requires iterative message passing over the resulting graph. Both are
shuffle-bound operations that Spark's execution model handles natively, with
spill-to-disk when a partition exceeds memory.

We deliberately used the \textbf{DataFrame} API rather than RDDs. An earlier
RDD implementation of this pipeline failed repeatedly with out-of-memory
errors, which we traced to the JVM-to-Python serialisation boundary: every row
crossing into a Python lambda is pickled, and the \texttt{stdout writer for
python} thread accumulated faster than it drained. Rewriting the same logic as
DataFrame expressions keeps the data and the computation inside the JVM, where
Catalyst can also push down projections and plan the joins. This change alone
turned a pipeline that could not complete into one that processes the full
corpus. It is the single most consequential engineering decision in the
project.

\textbf{Spark MLlib --- MinHashLSH.} Measuring review-text similarity within a
group is an all-pairs problem, and computing exact Jaccard similarity across
every pair of reviews in a component is not tractable at this scale. MinHash
locality-sensitive hashing approximates it by hashing shingled token sets into
bands, so that only candidate pairs that collide in a band are compared.
MLlib's \texttt{MinHashLSH} with \texttt{approxSimilarityJoin} gives this as a
distributed operation over a DataFrame of sparse vectors.

\textbf{GraphFrames 0.8.4.} Connected components on the account graph is an
iterative graph computation. GraphFrames provides a distributed implementation
backed by GraphX, invoked with \texttt{algorithm="graphx"} and a checkpoint
directory. Because GraphFrames is an external package resolved from the Spark
Packages repository at submit time, we also implemented a Pregel-style label
propagation fallback in pure DataFrame operations, selectable with
\texttt{USE\_GRAPHFRAMES=0}, so the pipeline runs in an offline or
dependency-restricted environment.

\textbf{AWS EMR on EC2, Amazon S3, and Apache Parquet.} EMR gave us a real
multi-node YARN cluster whose size we could change between otherwise identical
runs, which is exactly what the scale-out measurement requires. S3 holds both
the input corpus and the results, decoupling storage from the ephemeral
cluster, and Parquet is the output format because the results are read back
column-wise for analysis.

\subsection{Chosen Datasets}

We use the \textbf{Amazon Reviews'23} corpus released by the McAuley Lab at UC
San Diego (\url{https://amazon-reviews-2023.github.io/}). We selected four
categories that are large enough to be a genuine distributed workload while
spanning different product types and review cultures.

\begin{tabularx}{\textwidth}{@{}lXrr@{}}
  \toprule
  \textbf{Component} & \textbf{Contents} & \textbf{Compressed} & \textbf{Raw} \\
  \midrule
  User reviews &
    24{,}447{,}530 reviews across Video Games, CDs \& Vinyl, Arts Crafts \&
    Sewing, and Baby Products; 1996--2023 &
    3.71\,GB & 12.89\,GB \\
  Item metadata &
    1{,}858{,}398 products with store/brand, price, category hierarchy,
    catalogue average rating and rating count &
    1.05\,GB & 4.25\,GB \\
  \midrule
  \textbf{Total} & gzip-compressed JSONL (1\,GB $=10^9$ bytes)
    & \textbf{4.76\,GB} & \textbf{17.14\,GB} \\
  \bottomrule
\end{tabularx}

We verified these figures by decompressing and counting every file; the
per-file line counts sum to exactly the 24{,}447{,}530 reviews and
1{,}858{,}398 products the pipeline reports, which confirms that no records
are lost in ingest.

\begin{center}
\begin{tabular}{@{}lrrr@{}}
  \toprule
  \textbf{Category} & \textbf{Reviews} & \textbf{Products} &
  \textbf{Raw bytes} \\
  \midrule
  Arts, Crafts \& Sewing & 8{,}966{,}758 & 801{,}446 & 6.12\,GB \\
  Baby Products          & 6{,}028{,}884 & 217{,}724 & 3.64\,GB \\
  CDs \& Vinyl           & 4{,}827{,}273 & 701{,}959 & 4.24\,GB \\
  Video Games            & 4{,}624{,}615 & 137{,}269 & 3.12\,GB \\
  \midrule
  \textbf{Total}         & \textbf{24{,}447{,}530} & \textbf{1{,}858{,}398}
                         & \textbf{17.14\,GB} \\
  \bottomrule
\end{tabular}
\end{center}

\textbf{Structure.} Each review record is one JSON object per line with
\texttt{user\_id}, \texttt{asin}, \texttt{parent\_asin}, \texttt{rating},
\texttt{timestamp} (milliseconds since epoch), \texttt{text},
\texttt{helpful\_vote}, and \texttt{verified\_purchase}. Metadata records carry
\texttt{parent\_asin}, \texttt{main\_category}, \texttt{title},
\texttt{average\_rating}, \texttt{rating\_number}, \texttt{price}, and
\texttt{store}.

\textbf{Pre-processing pipeline.} We read both sources under
\emph{explicit} \texttt{StructType} schemas rather than letting Spark infer
them. This was not a cosmetic choice: schema inference requires a full extra
pass over the data, and on the metadata files it failed outright with a
\texttt{COLUMN\_ALREADY\_EXISTS} error, because the free-form
\texttt{details} object contains keys that differ only in case (for example
\texttt{Assembly Required} and \texttt{assembly required}). Declaring only the
eight review fields and six metadata fields we actually use avoids both
problems and reduces the bytes deserialised.

Reviews are then filtered for a non-null \texttt{user\_id},
\texttt{parent\_asin} and \texttt{timestamp}; metadata is deduplicated on
\texttt{parent\_asin}; and the two are joined. A review's
\texttt{rating\_deviation} is computed as its rating minus the product's
catalogue average. Reviews shorter than 30 characters are excluded
\emph{only} from the MinHash text analysis, not from the corpus, because a
short review still carries valid timing and rating evidence.

One property of the input shaped the whole performance story: the corpus
arrives as four gzip files, and \textbf{gzip is not splittable}. A gzip member
must be decompressed sequentially from its start, so Spark cannot divide one
file across tasks. The ingest stage is therefore capped at four parallel tasks
regardless of cluster size, which becomes important in
Section~\ref{sec:scaleout}.

% ======================================================================
\section{Detailed Component Description}
% ======================================================================

All logic lives in \texttt{AmazonReviewCoordinationDF.py} (612 lines). Every
tunable is an environment variable with a default, so the same file runs
locally and on EMR without modification.

\subsection{\texttt{load\_enriched} --- Ingest, Join, Feature Derivation}

Reads reviews and metadata under the explicit schemas, filters nulls,
deduplicates metadata on \texttt{parent\_asin}, and performs the enrichment
join. Derives \texttt{ts} (seconds), \texttt{year},
\texttt{review\_text} (null-coalesced), \texttt{helpful\_votes},
\texttt{verified}, \texttt{category}, \texttt{product\_title}, and
\texttt{rating\_deviation}. The join is a standard shuffle hash join on
\texttt{parent\_asin}; metadata is the smaller side.

\subsection{\texttt{add\_group\_key} --- Adaptive Time Bucketing}

Assigns each review a \texttt{group\_key} of the form
\texttt{parent\_asin \# bucket}. The bucket is
$\lfloor \texttt{ts} / (3600 \cdot 24) \rfloor$ by default. A fixed window is
wrong for popular products: a best-selling game can draw hundreds of unrelated
reviews in a day, and every pair of those reviewers would look like a
co-occurrence. The component therefore counts reviews per
(product, 24h-window) first, and for any window exceeding
\texttt{LARGE\_GROUP\_THRESHOLD} = 100 reviews, re-buckets that window at
\texttt{LARGE\_GROUP\_BUCKET\_HOURS} = 6 hours. This is a single
\texttt{groupBy}/\texttt{agg} followed by a broadcastable join and a
\texttt{when} expression, so the refinement costs one extra shuffle rather
than a per-product loop.

\subsection{\texttt{repeated\_pairs} --- Co-occurrence Self-Join}

The quadratic core. Distinct (\texttt{group\_key}, \texttt{user\_id}) rows are
self-joined on \texttt{group\_key} with the predicate
\texttt{user1 < user2}, which both removes self-pairs and canonicalises
ordering so each unordered pair appears once. The result is grouped by
(\texttt{user1}, \texttt{user2}) and counted, and only pairs with
\texttt{repeated\_groups} $\geq$ \texttt{MIN\_REPEATED\_GROUPS} = 3 survive.

Requiring three distinct co-occurrences, rather than one, is what makes the
output tractable and meaningful: two strangers can easily review the same
product on the same day once, but doing so three separate times is a strong
signal. Section~\ref{sec:sensitivity} quantifies how much this threshold
matters.

\subsection{\texttt{pair\_audit} and \texttt{drop\_same\_account\_pairs}}

A data-quality guard we added after inspecting early results. The corpus
contains user IDs that are the same base identifier with a numeric suffix
(for example \texttt{AHPG65LK...} and \texttt{AHPG65LK..\_1}). These produced
the highest co-occurrence counts in the graph --- one pair co-occurred 464
times --- but they almost certainly represent one underlying account, not two
coordinating ones. \texttt{pair\_audit} flags pairs whose IDs share a base
after suffix stripping, writes them to a CSV audit trail, and
\texttt{drop\_same\_account\_pairs} removes them before graph construction.
Keeping this as an explicit, auditable step rather than a silent filter means
the decision is visible and reversible.

\subsection{\texttt{connected\_components} and \texttt{label\_propagation}}

Builds a GraphFrame from the surviving pairs and runs
\texttt{connectedComponents(algorithm="graphx")} against a configured
checkpoint directory. Each resulting component is one candidate coordinated
group. \texttt{label\_propagation} is the fallback: it initialises each vertex
with its own ID, then iteratively propagates the minimum neighbour label for
up to 20 iterations, terminating early when no label changes. It is written
entirely in DataFrame operations with \texttt{localCheckpoint} to truncate the
lineage each round, which is necessary because iterative DataFrame
computations otherwise build a query plan that grows without bound.

\subsection{\texttt{text\_similarity} --- MinHash LSH}

Restricted to reviews belonging to graph components, capped at
\texttt{MAX\_REVIEWS\_PER\_COMPONENT} = 500 per component to bound the
all-pairs cost. Review text is lowercased, tokenised on non-word characters,
filtered to tokens of length $\geq 2$, and converted to a term-frequency
vector via \texttt{CountVectorizer}. \texttt{MinHashLSH} with
\texttt{MINHASH\_NUM\_HASH\_TABLES} = 5 and a fixed
\texttt{MINHASH\_SEED} = 42 produces hash signatures, and
\texttt{approxSimilarityJoin} returns candidate pairs within a Jaccard
distance corresponding to \texttt{TEXT\_SIMILARITY\_THRESHOLD} = 0.70.
Pairs are filtered to those whose two reviews sit in the same component and
belong to different accounts. The seed is fixed deliberately so that runs are
reproducible; MinHash is randomised, and without a fixed seed the similarity
counts would differ between the one-worker and two-worker runs and contaminate
the comparison.

\subsection{\texttt{component\_stats} and \texttt{score} --- Seven Signals}

\texttt{component\_stats} aggregates per component: user count, review count,
mean rating, mean absolute rating deviation from catalogue average, verified
rate, mean helpful votes, busiest single hour, edge count, mean repetition,
time concentration (busiest hour / total reviews), rating agreement
($1 - \mathrm{MAD}/4$), and graph density
($\mathrm{edges} / \binom{\mathrm{users}}{2}$).

\texttt{score} clamps each signal to $[0,1]$ and combines them:

\begin{center}
\begin{tabular}{@{}lr@{}}
  \toprule
  \textbf{Signal} & \textbf{Weight} \\
  \midrule
  Group size (normalised by 20 accounts)        & 0.18 \\
  Graph density                                  & 0.18 \\
  Repeated co-review frequency (normalised by 10) & 0.18 \\
  Text similarity (MinHash)                      & 0.14 \\
  Time concentration                             & 0.13 \\
  Deviation from catalogue rating (normalised by 2 stars) & 0.10 \\
  Rating agreement                               & 0.09 \\
  \midrule
  \textbf{Total}                                 & \textbf{1.00} \\
  \bottomrule
\end{tabular}
\end{center}

The weights are asserted to sum to 1.0 at import time. They encode a judgement
that structural evidence (size, density, repetition) is more trustworthy than
behavioural evidence (rating agreement), because the latter is easy to produce
by accident. We make no claim that these weights are optimal; they are a
documented, reproducible prior, and
Section~\ref{sec:sensitivity} reports how the structural thresholds behind
them affect the output.

\subsection{\texttt{suspicious\_reviews}}

Joins the enriched corpus back to the top-scoring components and writes the
individual reviews belonging to them, so that a human reviewer can inspect the
evidence behind a ranked group rather than only its score.

% ======================================================================
\section{Evaluation, Testing, and Execution}
% ======================================================================

\subsection{Evaluation Environment}

\textbf{Development and baseline (single node).} Linux workstation, PySpark
3.5 in local mode with \texttt{local[4]} and an 8\,GB driver heap. Used for the
full-corpus baseline, the sensitivity sweep, and the injection test.

\textbf{Cluster (AWS EMR on EC2).}

\begin{tabularx}{\textwidth}{@{}lX@{}}
  \toprule
  Release label            & \texttt{emr-7.14.0} \\
  Applications             & Hadoop 3.4.2, Spark 3.5.8-amzn-0, JDK 17 \\
  Cluster ID               & \texttt{j-05114253MUVDJMJSM0Y5} \\
  Primary node             & 1 $\times$ \texttt{m5.xlarge} (4 vCPU, 16\,GB) \\
  Core nodes               & 1 $\times$ \texttt{m5.xlarge} (Run A);
                             2 $\times$ \texttt{m5.xlarge} (Run B) \\
  Deploy mode              & YARN, client mode \\
  Driver memory            & 4\,GB \\
  Shuffle partitions       & 400 (identical in both runs) \\
  Storage                  & \texttt{s3://csc7740-amazon-review/} \\
  External package         & \texttt{graphframes:graphframes:0.8.4-spark3.5-s\_2.12} \\
  Cluster lifetime         & 7.78 hours, auto-terminated on a 1-hour idle policy \\
  \bottomrule
\end{tabularx}

Three EMR configuration problems were resolved during bring-up and are worth
recording because each produced a failure that looked unrelated to its cause.
All three traced to a single root: the application code set
\texttt{spark.hadoop.fs.defaultFS} to \texttt{file:///} in the
\texttt{SparkSession} builder. Builder options are applied on top of
\texttt{spark-submit} options when the context is created in client mode, so
this silently overrode the cluster's real \texttt{fs.defaultFS} and caused
every scheme-less or host-less path to resolve to the primary node's local
disk. The symptoms were: the ApplicationMaster on a core node being unable to
read a staging directory written to local disk; an
\texttt{Incomplete HDFS URI, no host} failure on EMR's default
\texttt{hdfs:///var/log/spark/apps} event-log path; and a
\texttt{FileNotFoundException} from \texttt{RawLocalFileSystem} when
\texttt{setCheckpointDir} resolved a bare path. The fix was to apply
\texttt{file:///} only when the master URL begins with \texttt{local}.

\subsection{Execution Instructions (Software Manual)}

\textbf{1. Prerequisites.} Python 3.9+, Java 17, and
\texttt{pip install -r requirements.txt} (PySpark 3.5, NumPy). GraphFrames is
resolved at submit time and needs network access; set
\texttt{USE\_GRAPHFRAMES=0} to use the built-in fallback instead.

\textbf{2. Obtain the data.}
\begin{lstlisting}
bash analysis/download_data.sh      # fetches 4 review + 4 metadata .jsonl.gz
\end{lstlisting}

\textbf{3. Run locally.}
\begin{lstlisting}
bash run_local.sh
# or, explicitly:
REVIEWS_PATH='data/reviews/*.jsonl.gz' \
METADATA_PATH='data/meta/*.jsonl.gz' \
OUTPUT_ROOT=output/df_full \
SPARK_MASTER='local[4]' DRIVER_MEMORY=8g SHUFFLE_PARTITIONS=200 \
python3 AmazonReviewCoordinationDF.py
\end{lstlisting}

\textbf{4. Run one category at a time}, useful on a memory-constrained
machine:
\begin{lstlisting}
bash run_all_categories.sh
\end{lstlisting}

\textbf{5. Deploy to AWS EMR.}
\begin{lstlisting}
# a) stage data and code
aws s3 cp data/reviews/ s3://<bucket>/data/reviews/ --recursive
aws s3 cp data/meta/    s3://<bucket>/data/meta/    --recursive
aws s3 cp AmazonReviewCoordinationDF.py s3://<bucket>/project/
aws s3 cp emr/run_emr.sh                s3://<bucket>/project/

# b) create the cluster (emr-7.14.0, Hadoop + Spark, m5.xlarge)
#    with a bootstrap action pointing at emr/bootstrap.sh

# c) submit the pipeline as a step
aws emr add-steps --cluster-id <id> --steps file://emr/step.json
\end{lstlisting}

\texttt{emr/run\_emr.sh} is the step entry point. It exports
\texttt{HADOOP\_CONF\_DIR}, queries HDFS for the real \texttt{fs.defaultFS},
pins the staging directory, event-log directory, SQL warehouse and checkpoint
directory to fully-qualified \texttt{hdfs://host:port} URIs, creates the
checkpoint directory up front, and then calls \texttt{spark-submit}.

\textbf{6. Reproduce the experiments.}
\begin{lstlisting}
python3 analysis/sensitivity.py   # window x threshold sweep
python3 analysis/injection.py     # synthetic ring recall test
python3 analysis/count_spark_operations.py   # Appendix 2 counts
\end{lstlisting}

\textbf{Environment variables.} Every tunable is overridable without editing
code:
\begin{lstlisting}
REVIEWS_PATH   METADATA_PATH    OUTPUT_ROOT      SPARK_MASTER
DRIVER_MEMORY  SHUFFLE_PARTITIONS                CHECKPOINT_DIR
USE_GRAPHFRAMES                 GRAPHFRAMES_PACKAGE
MINHASH_SEED   MAX_REVIEWS_PER_COMPONENT         WRITE_ENRICHED
\end{lstlisting}

\subsection{Results and Findings}

\subsubsection{Corpus-scale detection results}

The full pipeline run on all 24.4 million reviews produced identical results
on the single-node baseline and on both EMR cluster configurations:

\begin{center}
\begin{tabular}{@{}lr@{}}
  \toprule
  \textbf{Quantity} & \textbf{Value} \\
  \midrule
  Reviews ingested and enriched        & 24{,}447{,}530 \\
  Candidate product-time groups        & 20{,}640{,}171 \\
  Co-occurring account pairs ($\geq 3$ groups) & 1{,}051 \\
  Pairs flagged as same-account artefacts & 28 \\
  Pairs dropped before graph construction & 13 \\
  Edges in final graph                 & 1{,}038 \\
  Graph vertices (accounts)            & 1{,}294 \\
  Connected components                 & 443 \\
  \bottomrule
\end{tabular}
\end{center}

That three separate executions on different hardware and different cluster
topologies agree exactly on all eight figures is the strongest correctness
evidence we have: the pipeline is deterministic, and the fixed MinHash seed
holds.

\textbf{Interpretation.} 1{,}294 accounts out of millions, grouped into 443
components, is a small fraction --- and that is the honest headline. Most
components are account \emph{pairs}: of the top 20 by score, 17 have exactly
two members. The highest-scoring group (component
\texttt{1151051235330}, score 0.649) is two accounts that co-reviewed 11 times
with perfect text similarity across 64 reviews. The structurally largest
components are far bigger --- one has 223 accounts, 331 edges and 33{,}694
reviews --- but score lower (0.504) precisely because density falls as a
component grows, which is the intended behaviour of the density signal.

\subsubsection{Parameter sensitivity}
\label{sec:sensitivity}

Because the output size depends on two analyst choices --- the time window and
the co-occurrence threshold --- we swept both on the Video Games category
(4{,}624{,}615 reviews). The table reports surviving pairs, accounts, and
components, plus the largest component found.

\begin{center}
\begin{tabular}{@{}lrrrrrr@{}}
  \toprule
  \textbf{Window} & \textbf{Groups} & \textbf{Thresh.} & \textbf{Pairs} &
  \textbf{Accounts} & \textbf{Comps.} & \textbf{Largest} \\
  \midrule
  6h   & 3{,}997{,}723 & 2  & 1{,}532  & 2{,}376  & 1{,}051 & 30 \\
       &               & 3  & 58       & 95       & 44      & 4 \\
       &               & 5  & 1        & 2        & 1       & 2 \\
       &               & 10 & 0        & 0        & 0       & 0 \\
  \addlinespace
  12h  & 3{,}751{,}135 & 2  & 2{,}804  & 4{,}229  & 1{,}781 & 230 \\
       &               & 3  & 110      & 184      & 82      & 7 \\
       &               & 5  & 2        & 4        & 2       & 2 \\
       &               & 10 & 1        & 2        & 1       & 2 \\
  \addlinespace
  24h  & 3{,}428{,}395 & 2  & 4{,}989  & 7{,}037  & 2{,}784 & 675 \\
       &               & 3  & 224      & 338      & 129     & 34 \\
       &               & 5  & 10       & 16       & 7       & 4 \\
       &               & 10 & 1        & 2        & 1       & 2 \\
  \addlinespace
  72h  & 2{,}756{,}660 & 2  & 13{,}519 & 16{,}181 & 5{,}432 & 2{,}189 \\
       &               & 3  & 637      & 873      & 310     & 177 \\
       &               & 5  & 34       & 47       & 16      & 9 \\
       &               & 10 & 3        & 6        & 3       & 2 \\
  \addlinespace
  168h & 2{,}203{,}416 & 2  & 25{,}879 & 26{,}644 & 7{,}786 & 5{,}524 \\
       &               & 3  & 1{,}302  & 1{,}603  & 515     & 328 \\
       &               & 5  & 97       & 107      & 26      & 40 \\
       &               & 10 & 4        & 8        & 4       & 2 \\
  \bottomrule
\end{tabular}
\end{center}

Two findings matter. First, the threshold is far more powerful than the
window: at a fixed 24-hour window, moving from 2 to 3 co-occurrences cuts
surviving pairs by 95\% (4{,}989 $\rightarrow$ 224). Second, widening the
window inflates results superlinearly --- 168 hours at threshold 3 yields
5.8$\times$ the pairs of 24 hours --- and produces a 328-account component
that is almost certainly an artefact of popular products rather than
coordination. Our defaults (24h, threshold 3) sit at the knee of both curves.
Anyone reporting a count from a detector like this must publish its thresholds;
the same corpus supports answers spanning three orders of magnitude.

\subsubsection{Validation by synthetic injection}

Sensitivity tells us how the output moves, not whether it is right. To measure
recall we planted six synthetic coordination rings (125 reviews) into the
4{,}624{,}615-review Video Games corpus and re-ran detection at the default
settings.

\begin{center}
\begin{tabular}{@{}lrrrrl@{}}
  \toprule
  \textbf{Ring} & \textbf{Accts} & \textbf{Found} & \textbf{Recall} &
  \textbf{Foreign} & \textbf{Expected} \\
  \midrule
  R1 tight, 3 accounts, 5 products   & 3  & 3  & 100\% & 0 & detect \\
  R2 tight, 5 accounts, 3 products   & 5  & 5  & 100\% & 0 & detect \\
  R3 large, 10 accounts, 3 products  & 10 & 10 & 100\% & 0 & detect \\
  R4 below threshold, 2 products     & 5  & 0  & 0\%   & 0 & miss (2 $<$ 3) \\
  R5 staggered beyond window         & 5  & 0  & 0\%   & 0 & miss (outside 24h) \\
  R6 partial, 8 accounts, 6 products & 8  & 7  & 88\%  & 0 & partial \\
  \bottomrule
\end{tabular}
\end{center}

Every ring designed to be detectable was recovered completely, the two
designed to fall outside the stated operating envelope were correctly missed,
and --- importantly --- \textbf{zero foreign accounts} were pulled into any
planted ring's component. The detector does not merge unrelated accounts into
synthetic groups. The two deliberate misses are not defects but a measured
statement of the envelope: coordination spread across only two products, or
staggered beyond 24 hours, is invisible at the default settings and would
require lowering the threshold or widening the window, with the
false-positive cost quantified in the sensitivity table.

\subsubsection{Cluster scale-out: one worker versus two}
\label{sec:scaleout}

Both runs executed the identical code, commit, and settings; only the core
node count differed. Run A used one core node, Run B two.

\begin{center}
\begin{tabular}{@{}lrrr@{}}
  \toprule
  \textbf{Stage} & \textbf{1 worker} & \textbf{2 workers} & \textbf{Speedup} \\
  \midrule
  enrich + join                 & 80.9\,s   & 64.8\,s   & 1.25$\times$ \\
  candidate groups              & 178.2\,s  & 105.4\,s  & 1.69$\times$ \\
  repeated pairs (self-join)    & 40.8\,s   & 26.1\,s   & 1.56$\times$ \\
  pair audit                    & 1.9\,s    & 1.6\,s    & 1.19$\times$ \\
  connected components          & 6.4\,s    & 5.6\,s    & 1.14$\times$ \\
  MinHash text similarity       & 276.1\,s  & 175.8\,s  & 1.57$\times$ \\
  \textbf{component stats + score} & \textbf{6{,}381.2\,s} & \textbf{5{,}359.3\,s} & \textbf{1.19$\times$} \\
  write results                 & 91.1\,s   & 56.8\,s   & 1.60$\times$ \\
  \midrule
  \textbf{Total (wall clock)}   & \textbf{7{,}387\,s} & \textbf{6{,}046\,s} & \textbf{1.22$\times$} \\
  \bottomrule
\end{tabular}
\end{center}

Doubling the workers gave only a \textbf{1.22$\times$} overall speedup. The
reason is visible in the table: every stage except one scaled respectably
(1.14--1.69$\times$), but \texttt{component stats + score} consumed
\textbf{86.6\%} of Run A and improved by only 1.19$\times$, pinning the
overall figure to itself.

We profiled that stage rather than guessing. Three independent sources agree:

\begin{itemize}
  \item \textbf{CloudWatch.} The core node sat at a flat 27\% CPU --- almost
        exactly one of four vCPUs --- for over an hour, with EBS writes of
        roughly 20\,KB/s. The work was CPU-bound on a single thread at a time,
        not I/O-bound.
  \item \textbf{YARN metrics.} \texttt{ContainerAllocated} was 2 (one
        ApplicationMaster, one executor) with
        \texttt{YARNMemoryAvailablePercentage} at 1.04\% and
        \texttt{ContainerPending} around 12. Spark was requesting roughly six
        times the executors YARN could grant: Run A was resource-starved, not
        merely slow. After the resize, allocation rose to 3 containers and
        available memory to 39--50\%, confirming the second node was used.
  \item \textbf{Executor logs.} Stage 1482 completed 56 tasks between
        20:10:57 and 20:50:23, about 1.4 tasks per minute, each reading
        roughly 12\,MB of shuffle. Every fetch logged
        \texttt{Started 0 remote fetches}, i.e. all shuffle reads were local
        to the single node.
\end{itemize}

\textbf{Root cause.} In \texttt{component\_stats}, the intermediate
\texttt{member\_reviews = enriched.join(components, "user\_id")} is consumed
four times --- by the per-component aggregation, the busiest-hour
aggregation, the rating-agreement aggregation, and alongside edge statistics
--- and is never cached. Neither is \texttt{enriched}. Spark therefore
recomputes the entire ingest-and-join pipeline for each consumer, re-reading
and re-parsing 3.7\,GB of \emph{non-splittable} gzip every time. Because gzip
caps that re-read at four tasks, and Run A had only two executor cores, the
stage degenerated to roughly one core of useful work.

\textbf{Why the second worker helped so little.} Run B doubled task slots
from two to four, which should have halved a 139-task stage. It did not,
and the likely reason is the mirror image of the evidence above: with one
executor, Run A fetched every shuffle block locally; with two executors on
separate nodes, Run B must fetch roughly half of them across the network.
The added parallelism and the added network cost largely cancelled within
this stage. We state this as the probable mechanism rather than a confirmed
one --- distinguishing it conclusively would require comparing remote-fetch
counts in both runs' executor logs, which we have retained but not yet
analysed.

\textbf{The fix, and why we did not apply it mid-experiment.} Persisting
\texttt{member\_reviews} is a one-line change projected to cut total runtime
by 3--4$\times$. We deliberately left it out of these measurements so that
Run A and Run B remained comparable, and record it as the primary
optimisation for future work. Reporting the measured bottleneck with its
diagnosis is more useful than reporting a faster number with the cause
hidden.

\textbf{What this says about scalability.} The honest conclusion is that this
pipeline's scale-out is bounded not by the algorithm but by two implementation
properties: non-splittable input compression, which caps ingest parallelism at
the number of files, and an uncached intermediate, which multiplies that
capped work by four. Re-encoding the corpus as Parquet or splittable
block-compressed files, and caching the intermediate, would be prerequisites
to any meaningful study of scaling beyond two workers.

\textbf{Cost.} The entire cluster campaign --- bring-up, debugging, and both
full-corpus runs over a 7.78-hour cluster lifetime --- cost \$5.01:
\$4.39 in EC2/EMR compute, \$0.45 in S3 PUT requests, and \$0.17 in storage
and EBS.

\subsection{System Screenshots}

\begin{todobox}
Insert screenshots here before submission. The EMR cluster
\texttt{j-05114253MUVDJMJSM0Y5} is terminated but remains visible in the AWS
console for roughly two months, so all of these can still be captured.
Required: (1) EMR console cluster list showing the cluster and its
TERMINATED state; (2) the cluster's Summary tab showing release label
emr-7.14.0, applications, and instance counts; (3) the Hardware/Instance
groups tab showing the CORE group resized from 1 to 2 m5.xlarge; (4) the
Steps tab showing \texttt{one-worker-run-A} (s-0276858I3LVNHH4ZD8H) and
\texttt{two-worker-run-B} (s-09163011V7YC948GL1IZ) both COMPLETED with their
durations; (5) the S3 console showing
\texttt{s3://csc7740-amazon-review/results/one-worker/} and
\texttt{results/two-workers/} output prefixes; (6) the step stdout log
showing the stage timing table and PIPELINE COMPLETED SUCCESSFULLY block;
(7) CloudWatch graphs for ContainerAllocated and
YARNMemoryAvailablePercentage across both runs.
\end{todobox}

% ======================================================================
\section{Conclusion}
% ======================================================================

\textbf{What we accomplished.} We built and executed a distributed Spark
pipeline that processes 24.4 million Amazon reviews (17.1\,GB raw JSON) end to
end, joins them against 1.86 million product records, constructs an account
co-review graph, detects communities with GraphFrames, measures review-text
similarity with MinHash LSH, and ranks 443 candidate coordinated groups by a
seven-signal score. The pipeline ran to completion on a single node and on a
real two-configuration AWS EMR cluster, producing bit-identical analytical
results in all three executions.

We went beyond building the detector to measuring it. The sensitivity sweep
shows that the co-occurrence threshold moves the result by two orders of
magnitude and that reported counts are meaningless without stated parameters.
The injection test establishes 100\% recall on rings inside the operating
envelope with zero contamination of unrelated accounts. The scale-out
experiment quantifies a 1.22$\times$ speedup from doubling workers and
explains precisely why it is not 2$\times$.

\textbf{Lessons learned.} Three stand out. First, \emph{the API boundary is a
performance decision, not a style choice}: moving from RDDs with Python lambdas
to DataFrame expressions converted a pipeline that could not finish into one
that completes, because it removed the JVM-to-Python serialisation of every
row. Second, \emph{configuration defaults propagate in non-obvious ways}: a
single hard-coded \texttt{fs.defaultFS=file:///} intended for local runs
produced three unrelated-looking cluster failures, and we chased each symptom
separately before finding the common cause. Third, \emph{measure before
optimising, and profile before explaining}: our initial assumption about which
stage dominated was wrong, our first estimate of the remaining runtime was
wrong by a factor of three, and only executor-log and CloudWatch evidence
identified the real bottleneck as an uncached intermediate.

\textbf{Future improvements.} In priority order: (1) persist
\texttt{member\_reviews} and \texttt{enriched} to eliminate the four-fold
recomputation, the single highest-value change; (2) re-encode the corpus as
Parquet so ingest parallelism is not capped by gzip's non-splittability;
(3) confirm the remote-shuffle-fetch hypothesis and, if borne out, co-locate
partitions to reduce cross-node fetches; (4) extend to more categories and
measure whether coordination density varies by product type; (5) replace the
hand-chosen score weights with weights fitted against a labelled set, using
the injection harness to generate it; (6) add a streaming mode so new reviews
can be scored incrementally rather than by full recomputation.

\appendix

% ======================================================================
\section{Appendix 1: Source Code Summary and Engineering Contribution}
% ======================================================================

\begin{todobox}
Each team member must verify the AI-assistance column for their own files
before submission. The entries below describe the assistance actually used
during development as recorded in this project's working history.
\end{todobox}

{\small
\begin{longtable}{@{}p{1.6in}p{2.05in}p{2.05in}@{}}
  \toprule
  \textbf{File} & \textbf{Primary Functionality / Data Logic} &
  \textbf{AI Coding Assistant Usage Scope} \\
  \midrule
  \endhead
  \texttt{AmazonReview\allowbreak CoordinationDF.py} &
    Complete PySpark DataFrame pipeline: schema-explicit ingest, metadata
    join, adaptive time bucketing, co-occurrence self-join, GraphFrames
    connected components with label-propagation fallback, MinHash LSH text
    similarity, seven-signal scoring, same-account pair auditing &
    Extensive. Full DataFrame rewrite of an earlier RDD implementation was
    AI-assisted from a functional specification; schemas, scoring weights,
    pair-audit logic and the \texttt{WRITE\_ENRICHED} /
    \texttt{fs.defaultFS} fixes were AI-authored and human-reviewed \\
  \addlinespace
  \texttt{emr/run\_emr.sh} &
    EMR step entry point: resolves the cluster's real \texttt{fs.defaultFS},
    pins staging / event-log / warehouse / checkpoint URIs to fully-qualified
    \texttt{hdfs://host:port}, creates the checkpoint directory, invokes
    \texttt{spark-submit} with GraphFrames &
    AI-authored during cluster debugging; each configuration fix derived from
    a specific observed failure \\
  \addlinespace
  \texttt{emr/bootstrap.sh} &
    EMR bootstrap action installing NumPy, setuptools and GraphFrames on all
    nodes & AI-authored \\
  \addlinespace
  \texttt{analysis/}\allowbreak\texttt{sensitivity.py} &
    Time-window $\times$ co-occurrence-threshold parameter sweep, emitting
    pairs / accounts / components / largest-component per cell &
    AI-authored from a described experiment design \\
  \addlinespace
  \texttt{analysis/}\allowbreak\texttt{injection.py} &
    Synthetic coordination-ring injection and recall measurement, including
    foreign-account contamination check &
    AI-authored from a described experiment design \\
  \addlinespace
  \texttt{analysis/}\allowbreak\texttt{count\_spark\_}\allowbreak\texttt{operations.py} &
    AST-based static counter of Spark call sites for Appendix 2 &
    AI-authored \\
  \addlinespace
  \texttt{analysis/}\allowbreak\texttt{create\_sample\_}\allowbreak\texttt{dataset.py} &
    Builds the $<$10\,MB sample dataset for the deliverable package &
    \todo{confirm} \\
  \addlinespace
  \texttt{analysis/}\allowbreak\texttt{feasibility.py} &
    Early sizing analysis: review volume and group-size distribution &
    \todo{confirm} \\
  \addlinespace
  \texttt{analysis/}\allowbreak\texttt{pairest.py} &
    Co-occurrence pair-count estimation used to choose the threshold &
    \todo{confirm} \\
  \addlinespace
  \texttt{analysis/}\allowbreak\texttt{lockstep.py} &
    Exploratory lockstep-behaviour probe & \todo{confirm} \\
  \addlinespace
  \texttt{analysis/}\allowbreak\texttt{nearDup.py} &
    Exploratory near-duplicate review-text probe & \todo{confirm} \\
  \addlinespace
  \texttt{analysis/}\allowbreak\texttt{download\_data.sh} &
    Downloads the four review and four metadata archives from the McAuley Lab
    host & \todo{confirm} (host URL corrected with AI assistance) \\
  \addlinespace
  \texttt{run\_local.sh},\newline
  \texttt{run\_all\_}\allowbreak\texttt{categories.sh} &
    Local single-run and per-category sequential run drivers &
    \todo{confirm} \\
  \bottomrule
\end{longtable}
}

% ======================================================================
\section{Appendix 2: Distributed Operations Count}
% ======================================================================

Counts are static Spark call sites in the submitted Python source, obtained by
AST analysis (\texttt{analysis/count\_spark\_operations.py}). A call site
inside a loop is counted once. Local Python operations and Spark SQL
expression helpers such as \texttt{F.col()}, \texttt{F.avg()} and
\texttt{F.count()} are excluded, because they describe computation rather than
executing it. Counted operations include \texttt{agg},
\texttt{approxSimilarityJoin}, \texttt{cache}, \texttt{coalesce},
\texttt{collect}, \texttt{connectedComponents}, \texttt{count},
\texttt{distinct}, \texttt{drop}, \texttt{dropDuplicates}, \texttt{filter},
\texttt{fit}, \texttt{groupBy}, \texttt{join}, \texttt{json},
\texttt{limit}, \texttt{localCheckpoint}, \texttt{orderBy},
\texttt{parquet}, \texttt{persist}, \texttt{repartition}, \texttt{select},
\texttt{show}, \texttt{take}, \texttt{transform}, \texttt{unionByName},
\texttt{unpersist}, \texttt{where} and \texttt{withColumn}.

\begin{center}
\begin{tabular}{@{}lr@{}}
  \toprule
  \textbf{File} & \textbf{\# Distributed Operations} \\
  \midrule
  \texttt{AmazonReviewCoordinationDF.py} & 151 \\
  \texttt{analysis/sensitivity.py}       & 28 \\
  \texttt{analysis/injection.py}         & 22 \\
  \midrule
  \textbf{Total}                         & \textbf{201} \\
  \bottomrule
\end{tabular}
\end{center}

A per-call-site audit listing every operation with its line number is
included in the deliverable package as
\texttt{analysis/DISTRIBUTED\_OPERATIONS.md}.

\end{document}
